% Slajd 2 – problem jako koszt dla trzech stron i rynek, od którego zaczynamy
% (kryteria: związek z wyzwaniem 25 %, model biznesowy 20 %; RULES: Idea).
\begin{frame}{Kto dziś traci na braku informacji o~dostępności}
  \begin{columns}[T,onlytextwidth]
    \begin{column}{0.56\textwidth}
      \begin{block}{Trzy strony, jeden brak danych}
        \footnotesize
        \begin{itemize}
          \item \textbf{Gość} (wózek, wózek dziecięcy, senior): nie wie, czy wejdzie, więc nie rezerwuje albo jedzie
                na próżno. Etykieta „dostępne” bez źródła i~daty nic nie mówi.
          \item \textbf{Lokal:} traci klientów, których nigdy nie zobaczy, i~odpowiada na te same pytania o~schody
                i~toaletę mailem i~telefonem.
          \item \textbf{Miasto:} ma ustawowy obowiązek zapewniać dostępność, ale nie ma aktualnych danych
                ani zasobów na ręczne prowadzenie bazy.
        \end{itemize}
      \end{block}
      \begin{exampleblock}{Zakres prototypu}
        \footnotesize
        Osoby na wózkach i~rodzice z~wózkami dziecięcymi. Te same dane obsługują 8~profili potrzeb,
        od seniora po turystę medycznego, a~interfejs działa po polsku, angielsku, niemiecku i~ukraińsku.
      \end{exampleblock}
    \end{column}
    \begin{column}{0.42\textwidth}
      \textbf{\small Rynek, od którego zaczynamy}\\
      \muted{Kraków, katalog Accessly zbudowany z~OpenStreetMap}
      \vspace{1mm}

      \kpi{502}{noclegi}\hfill\kpi{2\,622}{gastronomia}\hfill\kpi{1\,117}{zdrowie}
      \vspace{2mm}

      \textbf{\small Popyt}\textsuperscript{*}
      \footnotesize
      \begin{itemize}
        \item ok.~14~mln odwiedzających Kraków rocznie
        \item 5,4~mln osób z~niepełnosprawnością w~Polsce; co piąty mieszkaniec ma 65+~lat
        \item rosnąca turystyka medyczna: pacjenci z~zagranicy pytają o~dostępność przed rezerwacją
      \end{itemize}
      \vspace{1mm}
      \muted{\textsuperscript{*} Dane zewnętrzne do zweryfikowania przed wystąpieniem: badania ruchu turystycznego
        (Małopolska Organizacja Turystyczna), GUS (NSP 2021, ludność).}
    \end{column}
  \end{columns}

  \note[item]{Zacząć od historii: Anna na wózku chce wyjść na kawę w~centrum. Mapa z~zeszłego roku mówi „dostępne”, a~dziś przy wejściu jest próg. Ona traci wieczór, kawiarnia traci klientkę, a~miasto nie wie, że próg istnieje.}
  \note[item]{Liczby po prawej z~katalogu to fakty z~naszej bazy (420 miejsc kultury też). Liczby z~gwiazdką to dane publiczne, które trzeba potwierdzić i~podać źródło, jeśli jury zapyta.}
  \note[item]{Grupa główna z~wyzwania: wózki i~wózki dziecięce. Pozostałe profile korzystają z~tych samych atrybutów, więc poszerzają rynek bez dodatkowych danych.}
\end{frame}
